\section*{الفصل \ref{ch:b1:weathering-clays} --- التجوية والطين}

\begin{solution}{exo:b1:weathering-clays:1}
\ce{K+}: $1/1.38 = 0.72$؛ \ce{Fe^2+}: $2/0.78 = 2.56$؛ كلاهما دون 3، أي كاتيونات قابلة للذوبان. \ce{Fe^3+}: $3/0.645 = 4.65$؛ \ce{Ti^4+}: $4/0.605 = 6.61$؛ كلاهما بين 3 و\num{10}، أي متحلّلات مائية. في تربة مشبعة بالماء، بلا أكسجين، يُطلَق الحديد ك\ce{Fe^2+}، فيبقى في المحلول ويحمله الماء؛ وفي تربة جيدة التصريف، يحوّله الأكسجين إلى \ce{Fe^3+}، فيترسّب فورًا كأكاسيد وهيدروكسيدات ويبقى.
\end{solution}

\begin{solution}{exo:b1:weathering-clays:2}
A يتمدّد بالغليكول ويتقلّص بالتسخين: \omterm{def:b1:weathering-clays:clay-mineral}{سميكتيت}. B عند \SI{7}{\angstrom} ويُتلف عند \SI{550}{\celsius}: \omterm{def:b1:weathering-clays:clay-mineral}{كاولينيت}. C، عند \SI{10}{\angstrom} مهما كانت المعاملة: \omterm{def:b1:weathering-clays:clay-mineral}{إيليت}. D، عند \SI{14}{\angstrom} دون تمدّد أو تقلّص: \omterm{def:b1:weathering-clays:clay-mineral}{كلوريت}.
\end{solution}

\begin{solution}{exo:b1:weathering-clays:3}
يعطي أورثوكلازان الألومنيومين في \omterm{def:b1:weathering-clays:clay-mineral}{كاولينيت} واحد؛ و$6 - 2 = 4$ سيليكونًا يغادر ك\ce{H4SiO4}؛ يغادر 2 \ce{K+}، توازنها 2 \ce{H+}؛ والأكسجين، $16 + w = 9 + 16$، يعطي $w = 9$:
\[
  \ce{2 KAlSi3O8 + 2 H+ + 9 H2O -> Al2Si2O5(OH)4 + 2 K+ + 4 H4SiO4}.
\]
الهيدروجين: $2 + 18 = 20 = 4 + 16$ يتحقّق. يغادر كل البوتاسيوم، و$4/6 =
2/3$ من السيليكون.
\end{solution}

\begin{solution}{exo:b1:weathering-clays:4}
(أ) الكتلة لا تتغيّر: $10 \times 999.87/916.72 = \qty{10.91}{cm^3}$.
(ب) بين نحو \qty{-4}{\celsius} و\qty{-15}{\celsius}. فقرب الصفر لا يستطيع الجليد أن يبني ضغطًا كافيًا لسدّ الشقّ؛ وفي صقيع أقسى لا يتحرّك الماء عبر المسامّ نحو الجليد، فيتوقّف الجليد عن النموّ. (ج) بحسب
\cref{prop:b1:weathering-clays:surface} مع $n = \qty{1}{m}/\qty{1}{cm} = 100$، يُضرب السطح في 100: من \qty{6}{m^2} إلى \qty{600}{m^2}.
\end{solution}

\begin{solution}{exo:b1:weathering-clays:5}
$S/m = 6/(\rho d)$: ل\qty{1}{cm}، $6/(2600\times0.01) =
\qty{0.231}{m^2.kg^{-1}}$؛ ل\qty{1}{mm}، \qty{2.31}{m^2.kg^{-1}}؛ ل\qty{1}{\micro m}، \qty{2308}{m^2.kg^{-1}}، أو \qty{2.3}{m^2.g^{-1}}. وتجري التفاعلات الكيميائية عند سطح الحبيبات، فيكبر معدّلها لكل كيلوغرام من الصخر مع السطح النوعي، عشرة آلاف مرة من السنتيمتر إلى الميكرومتر. \omterm{def:b1:weathering-clays:weathering}{التجوية الفيزيائية} تقطّع الصخر إلى قطع أصغر فأصغر، وتفتح شقوقًا يدخل منها الماء: فتضاعف السطح الذي تستطيع \omterm{def:b1:weathering-clays:weathering}{التجوية الكيميائية} أن تعمل فيه.
\end{solution}

\begin{solution}{exo:b1:weathering-clays:6}
يعادل طن الألبيت $\num{e6}/262.22 = \qty{3813.6}{mol}$. وبحسب \cref{prop:b1:weathering-clays:feldspars}، يعطي مولان واحدًا من \omterm{def:b1:weathering-clays:clay-mineral}{الكاولينيت}:
$\qty{1906.8}{mol} \times \qty{258.16}{g.mol^{-1}} = \qty{492.3}{kg}$. ويغادر اثنان من السيليكون الثلاثة:
$2\times3813.6 = \qty{7627.2}{mol}$ من \ce{SiO2}، \qty{458.2}{kg}. ويغادر كل الصوديوم: $3813.6\times22.99 =
\qty{87.7}{kg}$. ويبقى أقلّ من نصف كتلة \omterm{def:b1:rock-forming-minerals:groups}{الفلسبار} طينًا.
\end{solution}

\begin{solution}{exo:b1:weathering-clays:7}
بحسب \cref{prop:b1:weathering-clays:immobile}، يترك كيلوغرام من الجرانيت $0.30/0.60 = \qty{0.5}{kg}$ من السابرولايت. الكسر المفقود: الصوديوم $1 -
(0.35/3.5)(0.5) = 0.95$؛ السيليكون $1 - (60/72)(0.5) = 0.583$؛ الألومنيوم $1 -
(27/14)(0.5) = 0.036$. الألومنيوم محفوظ عمليًّا، والصوديوم قد زال تقريبًا كليًّا، ونحو \qty{60}{\percent} من السيليكون أيضًا، مع أن كثيرًا مما يبقى محجوز في الكوارتز: فالفلسبارات تتحوّل إلى \omterm{def:b1:weathering-clays:clay-mineral}{الكاولينيت}، \omterm{def:b1:weathering-clays:hydrolysis-degree}{التحلّل أحادي السيليكا}.
\end{solution}

\begin{solution}{exo:b1:weathering-clays:8}
البيروفيلايت: $x = y = 0$، شحنة 0. طبقة المسكوفيت: $x = 1$، شحنة $-1$. طبقة البيدلايت: $x = 0.5$، شحنة $-0.5$. طبقة المونتموريلونيت: $y = 0.33$، شحنة $-0.33$. وطبقتا \omterm{def:b1:weathering-clays:clay-mineral}{السميكتيت}، البيدلايت والمونتموريلونيت، تستطيعان أن تنتفخا: فشحنتهما الصغيرة يوازنها قليل من الكاتيونات بين الطبقات، فتمسك الطبقات بإحكام قليل، وتجذب ماءً حين تتميّه. والبيروفيلايت ليس فيه كاتيونات بين الطبقات لتتميّه، وشحنة المسكوفيت $-1$ يوازنها بوتاسيوم يقفل الطبقات.
\end{solution}

\begin{solution}{exo:b1:weathering-clays:9}
(أ) Mg: 2 و\num{2}؛ Si: 1 و\num{1}؛ C: 4 و\num{4}؛ O: $4 + 8 + 4 = 16$ و$12 + 4 =
16$؛ H: 8 و$4 + 4 = 8$؛ الشحنة: 0 و$4 - 4 = 0$. (ب) الكتلة المولية للفورستيريت هي $2\times24.305 + 28.085 + 4\times15.999 =
\qty{140.69}{g.mol^{-1}}$، فيكون الطن \qty{7108}{mol}. وهو يستهلك
$4\times7108 = \qty{28431}{mol}$ من \ce{CO2}، \qty{1251}{kg}؛ ومع اثنين من المغنيسيوم لكل وحدة صيغية، يخزّن \cref{prop:b1:weathering-clays:co2}
$2\times7108$ مولًا، \qty{626}{kg}، أي نصف المقدار المستهلك. (ج) طن من الكالسيت،
\qty{9991}{mol}، يستهلك \qty{9991}{mol} من \ce{CO2} (\qty{440}{kg})
عند ذوبانه، ويعيدها كلها عند ترسيب الكربونات من جديد: فلا يُخزَّن شيء.
\end{solution}

\begin{solution}{exo:b1:weathering-clays:10}
(أ) $3 < 3/r < 10$ يعطي $\qty{0.3}{\angstrom} < r < \qty{1.0}{\angstrom}$.
(ب) $3/0.27 = 11.1 > 10$: يتنقّل البورون كأكسيأنيون، هو البورات. (ج) $2/0.83 =
2.41$، كاتيون قابل للذوبان؛ و$4/0.53 = 7.55$، متحلّل مائي. وفي ماء فقير بالأكسجين يذوب المنغنيز ك\ce{Mn^2+}؛ وحيث يلتقي الينبوع بالهواء، يتأكسد إلى \ce{Mn^4+}، فيترسّب فورًا كأكسيد أسود على الصخر. (د) جهده الأيوني، $4/0.40 = 10.0$، يقع على الحدّ بين المتحلّلات المائية والأكسيأنيونات: فالسيليكون لا يُمسَك كالألومنيوم، ولا يُحمل بعيدًا كالكبريت، والكسر المغسول يتوقّف على كمية الماء المارّة، وهذا أساس الدرجات الثلاث للتحلّل المائي.
\end{solution}

\begin{solution}{exo:b1:weathering-clays:11}
(أ) بكتلة مولية للكربون \qty{12.011}{g.mol^{-1}}، فإن \num{11.7e12} إلى
\qty{13.6e12}{mol} تساوي من 0.141 إلى 0.163 مليار طن من الكربون في السنة؛ وبكتلة \qty{44.009}{g.mol^{-1}}، من 0.515 إلى 0.599 مليار طن من ثاني أكسيد الكربون. (ب) $4.08/13.6 = \qty{30}{\percent}$ إلى $4.08/11.7 = \qty{35}{\percent}$. (ج) يُخزَّن نصف ثاني أكسيد الكربون المستهلك حين يترسّب الكالسيوم والمغنيسيوم ككربونات: من 0.070 إلى 0.082 مليار طن من الكربون في السنة. والرقم الحقيقي أصغر، لأن جزءًا من البيكربونات يذهب مع الصوديوم والبوتاسيوم، اللذين لا يكوّنان كربونات في قاع البحر.
\end{solution}

\begin{solution}{exo:b1:weathering-clays:12}
Si/Al $= (w_{\ce{SiO2}}/60.08)/(2w_{\ce{Al2O3}}/101.96)$. A: $0.7748/0.7748 =
1.00$، نسبة \omterm{def:b1:weathering-clays:clay-mineral}{الكاولينيت}، \omterm{def:b1:weathering-clays:hydrolysis-degree}{التحلّل أحادي السيليكا}. B: $0.9987/0.4904 = 2.04$، \omterm{def:b1:rock-forming-minerals:mineral}{معدن} طيني 2:1، \omterm{def:b1:weathering-clays:hydrolysis-degree}{التحلّل ثنائي السيليكا}. C: $0.0333/1.2358 = 0.027$، بقي تقريبًا لا سيليكون، جيبسيت، \omterm{def:b1:weathering-clays:hydrolysis-degree}{التحلّل الألوميني}. ومن أقلّ الغسل إلى أشدّه، أي نحو المناخ الأحرّ والأرطب: B، A، C.
\end{solution}

\begin{solution}{pb:b1:weathering-clays:1}
\textbf{1.} \ce{2 NaAlSi3O8 + 2 H+ + 9 H2O -> Al2Si2O5(OH)4 + 2 Na+ + 4 H4SiO4}.

\textbf{2.} \ce{CaAl2Si2O8 + 2 H+ + H2O -> Al2Si2O5(OH)4 + Ca^2+}. لا يغادر السيليكون: فالأنورثيت يحتوي بالفعل على السيليكون الواحد لكل ألومنيوم في \omterm{def:b1:weathering-clays:clay-mineral}{الكاولينيت}.

\textbf{3.} جيبسيت واحد يأخذ ألومنيومًا واحدًا، أي نصف أنورثيت: فبالمضاعفة،
\ce{CaAl2Si2O8 + 2 H+ + 6 H2O -> 2 Al(OH)3 + Ca^2+ + 2 H4SiO4}. الأكسجين: $8 + 6 = 6 + 8$؛ الهيدروجين: $2 + 12 = 6 + 8$.

\textbf{4.} \ce{5 NaAlSi3O8 + 4 H+ + 16 H2O -> 2 Na_{0.5}Al2(Si_{3.5}Al_{0.5})O10(OH)2 + 4 Na+ + 8 H4SiO4}; % ce-unbalanced-ok: decimal subscripts the checker cannot read; balance tested in tests/bachelor-1/test_weathering_clays.py
الهيدروجين، $4 + 32 = 36$ على الطرف الأيسر و$2\times2 + 8\times4 = 36$ على
الأيمن، يتوازن.

\textbf{5.} كل \ce{H+} يأتي مع \ce{HCO3-} واحد من \ce{CO2} واحد و\ce{H2O}:
\ce{2 NaAlSi3O8 + 2 CO2 + 11 H2O -> Al2Si2O5(OH)4 + 2 Na+ + 2 HCO3- + 4 H4SiO4}
و\ce{CaAl2Si2O8 + 2 CO2 + 3 H2O -> Al2Si2O5(OH)4 + Ca^2+ + 2 HCO3-}.

\textbf{6.} الألبيت إلى البيدلايت هو \omterm{def:b1:weathering-clays:hydrolysis-degree}{التحلّل ثنائي السيليكا}؛ والألبيت والأنورثيت إلى \omterm{def:b1:weathering-clays:clay-mineral}{الكاولينيت} هما \omterm{def:b1:weathering-clays:hydrolysis-degree}{التحلّل أحادي السيليكا}؛ والأنورثيت إلى الجيبسيت هو \omterm{def:b1:weathering-clays:hydrolysis-degree}{التحلّل الألوميني}.

\textbf{7.} $\num{e6}/262.22 = \qty{3813.6}{mol}$ من الألبيت؛
$\num{e6}/278.20 = \qty{3594.5}{mol}$ من الأنورثيت.

\textbf{8.} الألبيت: $3813.6/2 = \qty{1906.8}{mol}$ من \omterm{def:b1:weathering-clays:clay-mineral}{الكاولينيت}،
\qty{492.3}{kg}. الأنورثيت: \qty{3594.5}{mol}، \qty{928.0}{kg}.

\textbf{9.} الألبيت: $2\times3813.6 = \qty{7627.2}{mol}$ من \ce{SiO2}،
\qty{458.2}{kg}، و\qty{3813.6}{mol} من الصوديوم، \qty{87.7}{kg}. الأنورثيت:
لا سيليكا، و\qty{3594.5}{mol} من الكالسيوم، \qty{144.1}{kg}.

\textbf{10.} الألبيت: $1/2.62 = \qty{0.382}{m^3}$؛ وكاولينيته،
$0.4923/2.6 = \qty{0.189}{m^3}$: فالمادة الصلبة تفقد \qty{50.4}{\percent} من
حجمها. الأنورثيت: $1/2.75 = \qty{0.364}{m^3}$؛ وكاولينيته،
$0.9280/2.6 = \qty{0.357}{m^3}$: خسارة لا تتعدّى \qty{1.8}{\percent}.

\textbf{11.} الفراغ المتروك يصير مسامّ، مملوءة بالماء والهواء: فالسابرولايت على جرانيت غنيّ بالألبيت خفيف ومسامي وهشّ، وحبيبات كوارتزه مفكّكة في القليل من الطين. أما الأنورثيت فيتحوّل إلى \omterm{def:b1:weathering-clays:clay-mineral}{كاولينيت} بحجم يقارب حجمه، فيحتفظ الغابرو المتجوّي بمعظم مادته الصلبة.

\textbf{12.} الصوديوم ($z/r = 0.98$) والكالسيوم (2.0) كاتيونات قابلة للذوبان وتغادر؛ والألومنيوم (5.6) متحلّل مائي ويبقى كله في \omterm{def:b1:weathering-clays:clay-mineral}{الكاولينيت}؛ والسيليكون (10.0)، على الحدّ، يغادر جزئيًّا من الألبيت، ولا يغادر البتّة من الأنورثيت، لأن \omterm{def:b1:weathering-clays:clay-mineral}{الكاولينيت} يستطيع أن يحفظ سيليكونه.

\textbf{13.} \ce{CO2} واحد لكل ألبيت: \qty{3813.6}{mol}، $3813.6 \times 44.01 =
\qty{167.8}{kg}$.

\textbf{14.} اثنان لكل أنورثيت: \qty{7189.0}{mol}، \qty{316.4}{kg}.

\textbf{15.} $7189.0 \times 61.02 = \qty{438.7}{kg}$ من البيكربونات.

\textbf{16.} تُرسّبه الكائنات ككربونات الكالسيوم،
\ce{Ca^2+ + 2 HCO3- -> CaCO3 + CO2 + H2O}: لكل كالسيوم، يعود أحد ال\ce{CO2} إلى الهواء، أي \qty{3594.5}{mol} أو \qty{158.2}{kg}، والآخر يُدفن في الكربونات.

\textbf{17.} الصوديوم لا يكوّن كربونات في قاع البحر: فبيكربونات لا يمكن تخزينه ككربونات صخرية، ويُحسب التخزين طويل المدى من الكالسيوم والمغنيسيوم وحدهما.

\textbf{18.} طن من الكالسيت يستهلك \qty{9991}{mol} من \ce{CO2} عند ذوبانه، ويعيدها كلها عند ترسيبه: لا يُخزَّن شيء، ولا يسحب ثاني أكسيد الكربون نهائيًّا إلا تجوية سيليكات الكالسيوم والمغنيسيوم.

\textbf{19.} $\qty{200}{m^3.s^{-1}} \times \qty{0.3}{mol.m^{-3}} =
\qty{60}{mol.s^{-1}}$؛ في سنة من \qty{3.156e7}{s}، \qty{1.893e9}{mol}، أو
$\num{1.893e9}/\num{50000} = \qty{37870}{mol}$ لكل كيلومتر مربع.

\textbf{20.} بيكربونتان لكل أنورثيت: \qty{9.467e8}{mol} في السنة،
$\num{9.467e8}\times\qty{278.20}{g} = \qty{263400}{t}$ في السنة، أو
\qty{5.27}{t} لكل كيلومتر مربع.

\textbf{21.} الصخر الطازج المُستهلك هو $5.27/0.20 = \qty{26.3}{t}$، أو
$\num{26340}/2700 = \qty{9.75}{m^3}$، لكل كيلومتر مربع وسنة: طبقة سمكها
\qty{9.75}{\micro m} في السنة، أي \qty{9.8}{m} لكل مليون سنة.

\textbf{22.} \ce{CO2} واحد مُخزَّن لكل كالسيوم: $\num{9.467e8} \times 44.01 =
\qty{4.17e10}{g}$، نحو \qty{41700}{t} في السنة.

\textbf{23.} إنه يحسب فقط المادة المحمولة في المحلول. \omterm{def:b1:weathering-clays:clay-mineral}{فالكاولينيت} والكوارتز اللذان تتركهما التجوية تُزالان أيضًا بالتعرية، كوحل ورمل في النهر، ولا يُحسبان؛ ولا يُحسب تجوية \omterm{def:b1:rock-forming-minerals:mineral}{المعادن} الأخرى في الصخر. وإن جاء جزء من البيكربونات من فلسبارات الصوديوم أو البوتاسيوم، التي تعطي بيكربونة واحدة لكل وحدة صيغية بدل اثنتين، لوجب أن يتجوّى صخر أكثر لتوفيرها.

\textbf{24.} بحسب \cref{prop:b1:weathering-clays:co2}، يُخزَّن مول واحد من \ce{CO2} لكل مول من الأنورثيت: $3594.5 \times \qty{44.01}{g} = \qty{158.19}{kg}$.
تجوية طن واحد من الأنورثيت تسحب نهائيًّا \textbf{158}~kg من ثاني أكسيد الكربون.
\end{solution}
